\section{Introduction}
\label{sec:rewrite-introduction}

Advances in artificial intelligence raise both hopes and fears. One hope is that increasingly capable AI could generate unprecedented economic prosperity. One fear is that workers who depend on labor income could be substituted by AI. In this paper, I show this hope and fear can coexist. In fact, I show that an permanent increase in wage growth can only happen when labor and AI are in fact substitutes. In this regime, labor's income share tends to zero despite  a higher wage growth, meaning that capital income and output per worker grow even faster. 

To show this, I extend the Ramsey--Cass--Koopmans model to include AI production services as an additional input in the production function. Competitive firms produce the final good with capital, effective labor, and AI production services. A monopolistic AI developer supplies these services using recursive AI self-improvement (RSI). Together, the elasticity of substitution between AI production services and labor and the maximum attainable level of AI efficiency determine whether human labor remains a long-run production bottleneck or whether AI accelerates economic growth.

%This fear, however, does not require real wages to decline: workers may receive a smaller share of a much larger economy even as their wages rise. Can unprecedented prosperity coexist with a vanishing labor income share? Under what conditions could the world economy follow such a path? How long could it take before the transition became visible, and what economic signals would reveal it? What would happen to wages and labor's income share along the way?

%This question is becoming more pressing as AI increasingly contributes to its own development. Recursive self-improvement (RSI) is a feedback process in which better AI helps develop still better AI. OpenAI \citep{openairesearchacceleration2026} and Anthropic \citep{favaroclark2026} report that AI systems already accelerate parts of their research. OpenAI's chief scientist already stated that the pace of progress ``could be sustained into recursive self-improvement'' \citep{pachocki2026}. Anthropic nevertheless cautions that, regarding AI systems that autonomously design and develop their successors, ``We are not there yet'' \citep{favaroclark2026}. These reports motivate studying the economic consequences of autonomous AI research without assuming when it will become possible.


The escape from the labor bottleneck arises from two feedback loops. A technological loop runs from RSI: greater AI efficiency produces more successful research from a given amount of research compute and thereby helps create still better AI. An economic loop runs from RSI to output and back to RSI: AI efficiency raises production and expands the resources that can finance research compute and therefore further AI improvements. Research compute, however, competes with the inference compute used in final-good production. %Household saving, capital accumulation, the allocation of compute between inference and research, and the developer's investment on AI efficiency improvement are determined jointly by the agent's decisions and the market clearing conditions.

I first characterize the labor-bottleneck regime. It arises when the elasticity of substitution between AI production services and labor is at most one, and can also arise when the elasticity exceeds one but the maximum attainable level of AI efficiency remains below an explicit threshold. In this regime, output per worker and the real wage eventually grow at the exogenous rate of labor-augmenting technological progress, and labor's income share converges to a positive value.

The outcome changes when the elasticity of substitution between AI production services and labor exceeds one and the maximum attainable level of AI efficiency lies above the threshold. An AI-dominated long-run regime then becomes possible. In this regime, capital and effective AI production services eventually grow faster than effective labor. The long-run growth rates of output per worker, the real wage and the interest rate exceed  their labor-bottleneck values. Because output per person grows faster than wages, labor's income share converges to zero. Still, AI-driven output growth raises wages despite the shrinking fraction of output paid to labor. %Within the AI-dominated regime, long-run growth and the interest rate increase with the maximum attainable level of AI efficiency.

I also examine the role of the upper bound on AI efficiency. With complementarity, long-run output-per-person growth cannot exceed the exogenous rate of labor-augmenting technological progress, even without an upper bound. At unit elasticity and without an upper bound, I derive conditions for balanced growth in which output per person and real wages grow faster than this exogenous rate. With substitution, AI profits may become unbounded, even with a finite-time horizon, ruling out the existence of an equilibrium.

I use numerical simulations to compare equilibrium paths between regimes, assuming that RSI suddenly becomes available in an economy that already uses AI but has not been having progress in AI efficiency. Growth initially rises even in
economies where labor remains a long-run production bottleneck. Moderate growth may persist for
a long time in the AI-dominated regime before output accelerates more sharply. In contrast, the growth rate in output per worker eventually returns to its long-run value in the labor-bottleneck regime. Depending on RSI productivity, the transition may look either like a slow take off or a sudden jump.

This paper has seven sections, including the present introduction. Section~\ref{sec:rewrite-literature} relates the paper to the literature. Section~\ref{sec:rewrite-model} presents the model and defines equilibrium. Section~\ref{sec:rewrite-regimes} establishes sufficient conditions for equilibrium with a finite upper bound on AI efficiency and characterizes the associated long-run growth and distributional outcomes. Section~\ref{sec:rewrite-uncapped} examines the three substitution regimes without an upper bound on AI efficiency. Section~\ref{sec:rewrite-quantitative} presents numerical simulations. Section~\ref{sec:rewrite-conclusion} concludes. The appendix contains the proofs and describes the algorithm used for the simulations.
